\documentclass[11pt]{article}
\usepackage[T1]{fontenc}
\usepackage[utf8]{inputenc}
\usepackage{lmodern}
\usepackage{amsmath,amssymb,amsthm}
\usepackage[margin=1in]{geometry}
\usepackage{hyperref}
\hypersetup{colorlinks=true,linkcolor=blue,citecolor=blue,urlcolor=blue}
\setlength{\parskip}{4pt}
\setlength{\parindent}{0pt}
\newtheorem{theorem}{Theorem}[section]
\newtheorem{lemma}[theorem]{Lemma}
\newtheorem{proposition}[theorem]{Proposition}
\newtheorem{corollary}[theorem]{Corollary}
\theoremstyle{remark}
\newtheorem{remark}[theorem]{Remark}
\newcommand{\Z}{\mathbb Z}
\newcommand{\Q}{\mathbb Q}
\newcommand{\C}{\mathbb C}
\newcommand{\Lcal}{\mathcal L}
\newcommand{\ord}{\operatorname{ord}}
\newcommand{\rank}{\operatorname{rank}}
\newcommand{\lcm}{\operatorname{lcm}}
\newcommand{\Res}{\operatorname{Res}}
\newcommand{\eps}{\varepsilon}
\title{A moving-denominator Hankel argument\protect\\for the irrationality of $\zeta_{1/p}(3)$}
\author{Open Problems research project}
\date{Research proof candidate --- 11 October 2026}

\begin{document}
\maketitle

\begin{center}
\begin{minipage}{0.94\textwidth}
\textbf{Review status.} This manuscript records a research proof candidate
assembled from internally audited, all-order arguments. It has not been
formalized in a proof assistant or independently validated outside this
research session. The theorem below is the claimed consequence of the
proof supplied here, not a statement of established publication status.
No part of the argument relies on finite determinant computations,
numerical sign patterns, or extrapolation from small sizes.
\end{minipage}
\end{center}

\begin{abstract}
For an integer $p\ge2$, put
$\xi_p=\sum_{n\ge1}n^2/(p^n-1)$.
We construct degree-$K$ rational polynomials from a Hankel matrix of a
second-order differential functional with denominator
$\prod_{a=1}^K(1-p^{-a}x)$.
The proposed irrationality argument has four ingredients: a real
determinant bound with cubic exponent $-K^3\log p/9$; a coefficientwise
factor $p^{K(K^2-1)/2}$; a cyclotomic denominator bound with leading
constant $385/(64\pi^2)$; and infinitely many nonzero determinants,
proved through rational interpolation of a nonrational Cauchy transform.
The inequality $385/(64\pi^2)<11/18$ then yields the contradiction to
rationality. We give proofs of all four ingredients, including the
polynomial-moment contribution to the cyclotomic denominator estimate.
\end{abstract}

\section{Statement, construction, and proof structure}

For $0<q<1$, we use the convention
\begin{equation}\label{eq:qzeta}
 \zeta_q(3)=\sum_{n=1}^{\infty}\frac{n^2q^n}{1-q^n}.
\end{equation}
The normalization matters: several inequivalent functions are called
$q$-zeta functions. In this manuscript $q=1/p$, where $p$ is an integer
at least $2$, and $\xi_p=\zeta_{1/p}(3)$.

\begin{theorem}[Candidate main theorem]\label{thm:main}
For every integer $p\ge2$, the real number
\[
 \xi_p=\sum_{n=1}^{\infty}\frac{n^2}{p^n-1}
\]
is irrational.
\end{theorem}

The literature search supporting this draft found no established proof
of the individual irrationality assertion in Theorem~\ref{thm:main}.
Jouhet and Mosaki prove that at least one of
$\zeta_q(3),\zeta_q(5),\zeta_q(7),\zeta_q(9)$ is irrational at reciprocal
integers~\cite{JM}. We do not regard that search as a proof of novelty.
Hankel determinants that are integer polynomials of growing degree and
have cubic exponential decay already occur in Zudilin's work on the
generalized $q$-logarithm~\cite{Zudilin2016}. Recent related work includes
accelerated identities for the same $q$-zeta value~\cite{Tauraso} and
explicit cyclotomic denominator calculations for other
$q$-series~\cite{KY}. The present candidate must be assessed through its
specific estimates, rather than through novelty of the general method.

Write $\theta=x\,d/dx$. For a function analytic on a neighborhood of
$[0,q]$, define
\begin{equation}\label{eq:functional}
 \Lcal(f)=\sum_{m=1}^{\infty}(\theta^2f)(q^m).
\end{equation}
This is well-defined: analyticity at $0$ implies
$\theta^2f(x)=O(x)$ there, and the evaluation points form a geometric
sequence. For $K\ge1$ let
\begin{equation}\label{eq:denom}
 D_K(x)=\prod_{a=1}^K(1-p^{-a}x),\qquad W_K(x)=D_K(x)^{-1}.
\end{equation}
All rational functions used below have poles outside $[0,q]$.

The elementary moments are
\begin{equation}\label{eq:poly-moments}
 \Lcal(1)=0,\qquad
 \Lcal(x^j)=\mu_j:=\frac{j^2}{p^j-1}\quad(j\ge1).
\end{equation}
Define
\begin{equation}\label{eq:Sa}
 S_a=\sum_{b=1}^a\frac{p^b(p^b+1)}{(p^b-1)^3}.
\end{equation}

\begin{lemma}[Pure simple-pole moments]\label{lem:purity}
For every $a\ge1$,
\[
 \Lcal\bigl((1-p^{-a}x)^{-1}\bigr)=\xi_p-S_a.
\]
\end{lemma}
\begin{proof}
The power-series identity
\[
 \sum_{n\ge1}n^2y^n=\frac{y(1+y)}{(1-y)^3}\qquad(|y|<1)
\]
and absolute convergence give
\begin{align*}
 \Lcal\bigl((1-q^a x)^{-1}\bigr)
 &=\sum_{m\ge1}\sum_{n\ge1}n^2q^{n(a+m)}\\
 &=\sum_{n\ge1}\frac{n^2q^n}{1-q^n}
   -\sum_{b=1}^a\sum_{n\ge1}n^2q^{bn}
 =\xi_p-S_a.
\end{align*}
The polynomial-moment formula follows by the same geometric summation.
\end{proof}

For an indeterminate $X$, define $\Lcal_X$ on polynomials and rational
functions with simple poles among $p,p^2,\ldots$ by
\begin{equation}\label{eq:formal-functional}
 \Lcal_X(1)=0,\quad \Lcal_X(x^j)=\mu_j\ (j\ge1),\quad
 \Lcal_X((1-p^{-a}x)^{-1})=X-S_a,
\end{equation}
extending linearly through the unique polynomial-plus-simple-pole
decomposition. Put
\begin{equation}\label{eq:Delta}
 \Delta_K(X)=\det\left[
 \Lcal_X\left(\frac{x^{i+j}}{D_K(x)}\right)
 \right]_{0\le i,j<K}\in\Q[X].
\end{equation}
Each entry is affine in $X$, so $\deg\Delta_K\le K$. Lemma~\ref{lem:purity}
shows that evaluation at $X=\xi_p$ recovers the actual real Gram
determinant of~\eqref{eq:functional}.

Let
\begin{equation}\label{eq:Edef}
 A_K=\frac{K(K+1)}2,\qquad
 e_K=\frac{K(K^2-1)}2=A_K(K-1).
\end{equation}
The proof constructs integers $C_K>0$, coprime to $p$, for which
\[
 P_K(X):=C_Kp^{-e_K}\Delta_K(X)\in\Z[X],
\]
and establishes
\begin{equation}\label{eq:overview-bound}
 |P_K(\xi_p)|\le
 \exp\left\{-\left(\frac{11}{18}-\frac{385}{64\pi^2}\right)
 (\log p)K^3+O_p(K^2\log K)\right\}.
\end{equation}
Infinitely many values $P_K(\xi_p)$ will be nonzero. If $\xi_p=a/b$ is
rational, $b^KP_K(a/b)$ is then a nonzero integer for infinitely many
$K$, contradicting~\eqref{eq:overview-bound}.

\section{An all-size real determinant bound}

\begin{proposition}\label{prop:analytic}
For $q=1/p\le1/2$ and $K\ge1$,
\begin{equation}\label{eq:analytic-exact}
 |\Delta_K(\xi_p)|\le
 \bigl[12(2K-1)^2\sqrt K\bigr]^Kq^{b_K},
 \qquad
 b_K=\frac{K^3+3K^2-c_{K\bmod3}}9,
\end{equation}
where $(c_0,c_1,c_2)=(0,4,2)$. In particular,
\[
 \log|\Delta_K(\xi_p)|\le-\frac{\log p}{9}K^3+O_p(K^2),
\]
with the logarithmic inequality understood in the extended sense when
the determinant vanishes.
\end{proposition}

\begin{proof}
Introduce monic polynomials
\[
 B_j(x)=\prod_{t=1}^j(x-q^{\lceil t/3\rceil}),\quad B_0=1,
 \quad R_j=\lceil j/3\rceil,\quad
 a_j=\sum_{t=1}^j\lceil t/3\rceil.
\]
Replacing $1,x,\ldots,x^{K-1}$ by $B_0,\ldots,B_{K-1}$ is a triangular
basis change with determinant one. The Gram determinant is unchanged.

If $m<R_j$, the polynomial $B_j$ has a root of multiplicity three at
$q^m$. Hence $(\theta^uB_j)(q^m)=0$ for $0\le u\le2$. If $m\ge R_j$,
write $x=q^m$. Every root $\rho$ of $B_j$ satisfies $x\le\rho$, so
$|x-\rho|\le\rho$. Moreover, $\theta^u(x-\rho)=x$ for $u\ge1$.
Applying the product rule once or twice therefore gives
\begin{equation}\label{eq:jet-bounds}
 |B_j(x)|\le q^{a_j},\qquad
 |\theta^uB_j(x)|\le j^uq^{a_j}\frac{x}{q^{R_j}}\quad(u=1,2).
\end{equation}
For completeness, there are $j^u$ labeled choices of factors under $u$
successive differentiations. Each differentiated factor replaces a
bound $\rho$ by $x$; all ratios $x/\rho$ are at most one, and at least
one is at most $x/q^{R_j}$. For $j=0$ the positive derivatives vanish.

We next bound $W_K$ and its two derivatives uniformly in $K$.
For $0\le x\le q\le1/2$, define
\[
 \alpha(x)=\theta\log W_K(x)
   =\sum_{a=1}^K\frac{q^ax}{1-q^ax},\qquad
 \beta(x)=\theta^2\log W_K(x)
   =\sum_{a=1}^K\frac{q^ax}{(1-q^ax)^2}.
\]
Since $1-q^ax\ge3/4$ and $\sum_{a\ge1}q^a\le1$,
\[
 \log W_K(x)\le\alpha(x)\le\tfrac43x\le\tfrac23,
 \qquad \beta(x)\le\tfrac{16}{9}x.
\]
Using $\theta W_K=W_K\alpha$, $\theta^2W_K=W_K(\alpha^2+\beta)$,
and $x^2\le x/2$, we obtain
\begin{equation}\label{eq:weight-bounds}
 |W_K|\le e^{2/3},\qquad
 |\theta W_K|\le\tfrac43e^{2/3}x,\qquad
 |\theta^2W_K|\le\tfrac83e^{2/3}x.
\end{equation}

Fix $i,j<K$ and put $R=\max(R_i,R_j)$. The product
$\theta^2(B_iB_jW_K)$ vanishes at every $q^m$ with $m<R$. At a surviving
node, put $h=i+j$. Equations~\eqref{eq:jet-bounds}--\eqref{eq:weight-bounds}
and the product rule imply
\begin{align*}
 |\theta^2(B_iB_jW_K)(q^m)|
 &\le e^{2/3}\left(h^2+\tfrac83h+\tfrac83\right)
       q^{a_i+a_j}q^{m-R}\\
 &\le\tfrac83e^{2/3}(h+1)^2q^{a_i+a_j}q^{m-R}.
\end{align*}
Here every term contains a derivative, which provides at least one
factor $x/q^R$; any additional such factors are at most one.
Summing over $m\ge\max(1,R)$ costs at most $(1-q)^{-1}\le2$.
Since $(16/3)e^{2/3}<12$,
\begin{equation}\label{eq:entry-bound}
 |\Lcal(B_iB_jW_K)|\le12(i+j+1)^2q^{a_i+a_j}.
\end{equation}

Extract $q^{a_i}$ from row $i$ and $q^{a_j}$ from column $j$.
Hadamard's inequality applied to the remaining matrix gives
\[
 |\Delta_K(\xi_p)|\le
 q^{2\sum_{j=0}^{K-1}a_j}\bigl[12(2K-1)^2\sqrt K\bigr]^K.
\]
If $j=3u+r$, $0\le r<3$, direct summation gives
$a_j=3u(u+1)/2+r(u+1)$. Summing this expression over the three residue
classes gives $2\sum_{j<K}a_j=b_K$ in~\eqref{eq:analytic-exact}.
Finally the logarithm of the prefactor is $O(K\log K)$.
\end{proof}

\section{The coefficientwise factor at primes dividing $p$}

\begin{proposition}\label{prop:content}
For every prime $\ell\mid p$, put $r=v_\ell(p)$. Every coefficient of
$\Delta_K(X)$ has $\ell$-adic valuation at least $re_K$. Furthermore,
\begin{equation}\label{eq:leading}
 [X^K]\Delta_K(X)=(-1)^{K(K-1)/2}p^{e_K}.
\end{equation}
In particular $\Delta_K$ has degree exactly $K$ and
$p^{-e_K}\Delta_K\in\Z_\ell[X]$ for each prime $\ell\mid p$.
\end{proposition}

\begin{proof}
Set $r_a=p^a$ and
\[
 F_K(x)=\prod_{a=1}^K(x-r_a),\qquad
 D_K(x)=(-1)^Kp^{-A_K}F_K(x).
\]
Let $h_t=h_t(r_1,\ldots,r_K)$ be the complete homogeneous symmetric
polynomial, with $h_0=1$ and $h_t=0$ for $t<0$. It follows directly
from its monomial expansion that $h_t\in\Z$ and
$v_\ell(h_t)\ge rt$ for $t\ge0$.

Work in $\Q_\ell$. Introduce
\[
 g(y)=\frac{y(1+y)}{(1-y)^3},\quad
 Z=\sum_{n\ge1}g(p^n),\quad
 c_j=\frac{j^2p^j}{1-p^j},\quad
 H(t)=\sum_{j\ge1}c_jt^j=\sum_{n\ge1}g(p^nt).
\]
All required series converge $\ell$-adically; in particular
$v_\ell(c_j)\ge rj$. The identity $g(y^{-1})=-g(y)$ gives
\begin{equation}\label{eq:local-tail}
 S_a=-Z+H(r_a).
\end{equation}
To check it explicitly, the definition of $S_a$ equals
$-\sum_{n=1}^ag(p^n)$, while $H(r_a)=\sum_{n>a}g(p^n)$.

Put
\[
 M_k(X)=\Lcal_X\left(\frac{x^k}{F_K(x)}\right).
\]
The quotient in the division of $x^k$ by $F_K$ is zero for $k<K$;
for $k\ge K$ it is
$\sum_{j=0}^{k-K}h_{k-K-j}x^j$.
The partial-fraction coefficient at $x=r_a$ is
$r_a^k/F_K'(r_a)$, and
\[
 \Lcal_X((x-r_a)^{-1})=-\frac{X-S_a}{r_a}.
\]
The divided-difference identities needed below are
\begin{equation}\label{eq:dd}
 \sum_{a=1}^K\frac{r_a^m}{F_K'(r_a)}=h_{m-K+1}\quad(m\ge0),
 \qquad
 \sum_{a=1}^K\frac1{r_aF_K'(r_a)}=-\frac1{F_K(0)}.
\end{equation}
The first identity follows by expanding
$1/F_K(x)=\sum_a(F_K'(r_a)(x-r_a))^{-1}$ at infinity;
the second is its evaluation at $x=0$.

Write $U=X+Z$ and define
\begin{equation}\label{eq:Rk}
 R_k=\sum_{j\ge1}c_jh_{k+j-K}
       +\sum_{j=1}^{k-K}\mu_jh_{k-K-j},
\end{equation}
where an upper limit below one means an empty sum.
Substituting~\eqref{eq:local-tail} into the partial fractions and then
using~\eqref{eq:dd} gives the exact identities
\begin{equation}\label{eq:Mk}
 M_0(X)=\frac{U}{F_K(0)}+R_0,\qquad
 M_k(X)=-Uh_{k-K}+R_k\quad(k\ge1).
\end{equation}
The infinite sum is convergent: its nonzero terms have valuation at
least $rj+r(k+j-K)$ once $k+j-K\ge0$.
For $k=0$, the divided difference is applied to $H(t)/t$, whose
power-series exponents are nonnegative, so no additional term is
missing. All $R_k$ are $\ell$-integral, since the $h_t$ and $c_j$ are
integral and $\mu_j=j^2/(p^j-1)$ is integral at $\ell$.

Let $T=[M_{i+j}(X)]_{0\le i,j<K}$. Every entry except possibly $(0,0)$
belongs to $\Z_\ell[X]$. Since
$v_\ell(F_K(0))=rA_K$, multiplying row zero by $p^{A_K}$ makes $T$
integral over $\Z_\ell[X]$. Therefore every coefficient of $\det T$
has valuation at least $-rA_K$. The scalar relationship is
\begin{equation}\label{eq:Delta-T}
 \Delta_K(X)=(-1)^{K^2}p^{A_KK}\det T.
\end{equation}
This proves the coefficient lower bound $rA_K(K-1)=re_K$.

For the leading coefficient, the matrix of coefficients of $X$ in $T$
has $(0,0)$ entry $1/F_K(0)$ and zero entries elsewhere in row or
column zero. Its remaining block has entries $-h_{i+j-K}$ for
$1\le i,j<K$: entries with $i+j<K$ vanish, and entries with $i+j=K$
equal $-1$. Reversing the columns makes this block triangular.
Its determinant is
$(-1)^{(K-1)(K-2)/2}(-1)^{K-1}$.
Since $F_K(0)=(-1)^Kp^{A_K}$, substitution
in~\eqref{eq:Delta-T} gives~\eqref{eq:leading}, including $K=1$.
\end{proof}

\section{A cyclotomic denominator estimate}

This section works first with an indeterminate $z$ in place of $p$.
Let $\Delta_{K,z}(X)$ be defined by~\eqref{eq:formal-functional}
and~\eqref{eq:Delta}, replacing $p$ everywhere by $z$. Write
$\Phi_d(z)$ for the $d$th cyclotomic polynomial. Define
\begin{equation}\label{eq:cd}
 c_{K,d}=\begin{cases}
 \left\lceil K^2/(4d)+2K+17d/16\right\rceil,&1\le d\le K/2,\\
 7(K-d)+3,&K/2<d\le K,
 \end{cases}
\end{equation}
and set
\begin{equation}\label{eq:CK}
 C_K=\prod_{d=1}^K\Phi_d(p)^{c_{K,d}}.
\end{equation}

\begin{proposition}\label{prop:denominators}
One has
\begin{equation}\label{eq:laurent-integral}
 \left(\prod_{d=1}^K\Phi_d(z)^{c_{K,d}}\right)\Delta_{K,z}(X)
       \in\Z[z,z^{-1}][X].
\end{equation}
For every fixed integer $p\ge2$, $C_K$ is a positive integer coprime
to $p$, and
\begin{equation}\label{eq:denom-growth}
 \log C_K\le\frac{385}{64\pi^2}(\log p)K^3
                   +O_p(K^2\log K).
\end{equation}
\end{proposition}

We separate local pole bounds from their global summation. All local
estimates apply coefficientwise in $X$.

\subsection{The cluster basis at a root of unity}

Fix $d\le K$, a primitive $d$th root $\omega$, and the uniformizer
$\tau=z-\omega$. Let $v=\ord_\tau$ and work over $\C((\tau))$.
The variable $X$ remains formal; the valuation of a polynomial in $X$
means the minimum valuation of its coefficients. Set
\[
 F(x)=\prod_{a=1}^K(x-z^a),\qquad
 D(x)=\prod_{a=1}^K(1-z^{-a}x).
\]
Partition the indices into residue classes modulo $d$:
\[
 I_r=\{a:1\le a\le K,\ a\equiv r\pmod d\},\quad
 m_r=|I_r|,\quad F_r(x)=\prod_{a\in I_r}(x-z^a).
\]
For each nonempty cluster use the polynomials
\begin{equation}\label{eq:cluster-basis}
 G_{r,j}(x)=\frac{F(x)}{F_r(x)}(x-\omega^r)^j,
 \qquad 0\le j<m_r.
\end{equation}
There are $K$ such polynomials, all of degree below $K$, and their
change-of-basis determinant is a local unit. Indeed, at $\tau=0$ the
polynomials from cluster $r$ vanish to full multiplicity at all other
clusters. Modulo $(x-\omega^r)^{m_r}$ they are the ordinary powers
multiplied by a unit. The Chinese remainder theorem therefore proves
that their reductions form a basis. The Gram determinant in this
basis has the same local valuation as $\Delta_{K,z}$.

Divide $G_{r,j}G_{s,k}$ by $D$ and decompose its proper part into
simple poles. The polynomial quotient has degree at most $K-2$ and
coefficients in $\C[[\tau]]$; the leading coefficient of $D$ is a unit.
Applying the polynomial moments to this quotient produces at most a
simple pole, because each $j^2/(z^j-1)$ has at most a simple pole.

If $r\ne s$, $F$ divides $G_{r,j}G_{s,k}$, so the proper part is zero.
If $r=s$, the coefficient of a simple pole at $x=z^a$, $a\in I_r$,
has valuation at least
$j+k-(m_r-1)$. To see this, the numerator contributes
$(z^a-\omega^r)^{j+k}$, while $F'(z^a)$ has precisely $m_r-1$
vanishing factors, each of order one. All other factors are units.
The pole moment $X-S_a(z)$ has order at most three. Thus the image of
the proper part has valuation at least
\begin{equation}\label{eq:pole-local}
 j+k-m_r-2.
\end{equation}

\subsection{A general bound for small cyclotomic orders}

In a cluster of size $m$ assign the weights
\[
 w_j=\max\{1/2,(m+2)/2-j\},\qquad 0\le j<m.
\]
Multiply row and column $j$ by $\tau^{w_j}$. Half-integer powers can
be interpreted in $\C[[\sqrt\tau]]$, retaining the normalization
$v(\tau)=1$. Every scaled polynomial-part entry is integral because
the sum of any two weights is at least one. Every proper-part entry
is integral by~\eqref{eq:pole-local} and
$w_j+w_k\ge m+2-j-k$.
It follows that
\begin{equation}\label{eq:cluster-cost}
 -v(\Delta_{K,z})\le\sum_r f(m_r),\qquad
 f(m)=\sum_{j=0}^{m-1}\max\{1,m+2-2j\}.
\end{equation}
Elementary summation gives
\[
 f(m)=\begin{cases}
 m^2/4+2m+1,&m\text{ even},\\
 m^2/4+2m+3/4,&m\text{ odd}.
 \end{cases}
\]
In particular $f(m)\le m^2/4+2m+1$. Write $K=td+s$, $0\le s<d$.
There are $s$ clusters of size $t+1$ and $d-s$ of size $t$, whence
\[
 \sum_r m_r^2=\frac{K^2}{d}+s(1-s/d)
 \le\frac{K^2}{d}+\frac d4.
\]
Combining this with~\eqref{eq:cluster-cost} yields
\begin{equation}\label{eq:low-order}
 -v(\Delta_{K,z})\le\frac{K^2}{4d}+2K+\frac{17d}{16}.
\end{equation}
This proves the required bound in the first range of~\eqref{eq:cd}.

\subsection{The high-order range and the polynomial-moment correction}

Now assume $K/2<d\le K$, and put $t=K-d$. The clusters have sizes one
or two, with precisely $t$ double clusters. Use weights $(2,1)$ on
each double cluster, weight $3/2$ on the singleton with index $a=d$,
and weight zero on all other singletons. The last singletons have
indices $a<d$, so their $S_a(z)$ are analytic at $z=\omega$ and their
proper-part entries have no pole. The double clusters are made
integral by~\eqref{eq:pole-local}; the $a=d$ singleton has at most a
third-order pole. Cross-cluster proper parts vanish. These scalings
therefore clear the entire proper part, at total determinant cost
\begin{equation}\label{eq:high-scale}
 6t+3.
\end{equation}

The polynomial quotient still needs to be included. Its degree is at
most $K-2<2d$, so its only possibly singular nonconstant moment is
\[
 \mu_d(z)=\frac{d^2}{z^d-1}.
\]
It can occur when $K-d\ge2$; the implication $d>K/2\Rightarrow K-2<d$
is false and is not used here.

For polynomials $f,g$ of degree below $K$, let
\[
 B_d(f,g)=[x^d]\,(fg/D)_+,
\]
where $(\cdot)_+$ denotes polynomial quotient. If $\deg f\le d$, then
$\deg(fg)\le K+d-1$, and its quotient by $D$ has degree below $d$.
Thus $B_d(f,g)=0$ for every $g$. The radical of this bilinear form
contains the polynomials of degree at most $d$, and consequently
\begin{equation}\label{eq:rank}
 \rank B_d\le r_d:=\max\{0,K-d-1\}=\max\{0,t-1\}.
\end{equation}
For $d\ge K-1$ the form is identically zero, as also stated by this
formula. This is an exact rank bound over the fraction field, not only
a rank bound after reduction modulo $\tau$.

The coefficient matrix of $B_d$ is integral at $\tau=0$. The cluster
basis change and all row and column weights preserve its rank and
integrality. Following the scalings in~\eqref{eq:high-scale}, the full
Gram matrix therefore has the form
\[
 A(\tau)+\tau^{-1}B(\tau),
\]
where $A,B$ have coefficients in $\C[[\sqrt\tau]][X]$ and
$\rank B\le r_d$. The analytic unit
$\tau d^2/(z^d-1)$ has been absorbed into $B$.
In the determinant expansion by columns, every term choosing more
than $r_d$ columns from $B$ vanishes by linear dependence. Each
remaining term has valuation at least $-r_d$. Removing the row and
column scalings proves
\begin{equation}\label{eq:high-order}
 -v(\Delta_{K,z})\le6t+3+\max(0,t-1)\le7(K-d)+3.
\end{equation}
This proves the high-order part of~\eqref{eq:cd} without discarding
the polynomial moments.

\subsection{Returning to integral polynomials}

The original partial fractions and polynomial division show directly
that every coefficient of $\Delta_{K,z}$ belongs to
\begin{equation}\label{eq:localization}
 \Z[z,z^{-1},(z^n-1)^{-1}:1\le n\le K].
\end{equation}
Indeed the differences between distinct pole positions are
$z^a-z^b$, each a monomial times $z^{|a-b|}-1$; the moments introduce
only the displayed types of factors. In this ring all denominators
other than powers of $z$ factor into the monic polynomials $\Phi_d$
with $d\le K$. The local estimates at their primitive roots cancel
every such denominator. Gauss's lemma, or divisibility by monic
polynomials in $\Z[z]$, then gives~\eqref{eq:laurent-integral}. There
are no additional rational constant denominators in~\eqref{eq:localization}.

At $z=p\ge2$ all $\Phi_d(p)$ are positive integers. If $\ell\mid p$,
then $\Phi_d(p)\equiv\Phi_d(0)\equiv\pm1\pmod\ell$, so $C_K$ is
coprime to $p$.

\subsection{Summation of the local costs}

We use the elementary estimates
\begin{align}
 \sum_{d\le M}\frac{\varphi(d)}d
   &=\frac6{\pi^2}M+O(\log M),\label{eq:tot1}\\
 \sum_{d\le M}\varphi(d)
   &=\frac3{\pi^2}M^2+O(M\log M),\label{eq:tot2}\\
 \sum_{d\le M}d\varphi(d)
   &=\frac2{\pi^2}M^3+O(M^2\log M).\label{eq:tot3}
\end{align}
To recall their elementary derivation, use
$\varphi(n)/n=\sum_{a\mid n}\mu(a)/a$, where here $\mu$ is the
M\"obius function. Summing and using
$\sum_{a\ge1}\mu(a)/a^2=1/\zeta(2)=6/\pi^2$ proves~\eqref{eq:tot1};
the floor-function errors are bounded by $\sum_{a\le M}1/a$ and the
tail of the absolutely convergent series is $O(1/M)$.
Partial summation gives~\eqref{eq:tot2} and~\eqref{eq:tot3}.

The ceiling in the low-order range contributes only
$O(\sum_{d\le K}\varphi(d))=O(K^2)$. Applying these estimates with
$M=K/2$ gives
\begin{align*}
 \sum_{d\le K/2}c_{K,d}\varphi(d)
 &\le\frac{K^2}{4}\sum_{d\le K/2}\frac{\varphi(d)}d
   +2K\sum_{d\le K/2}\varphi(d)
   +\frac{17}{16}\sum_{d\le K/2}d\varphi(d)+O(K^2)\\
 &=\frac{161}{64\pi^2}K^3+O(K^2\log K).
\end{align*}
For the high-order range,
\begin{align*}
 \sum_{K/2<d\le K}c_{K,d}\varphi(d)
 &=7K\sum_{K/2<d\le K}\varphi(d)
     -7\sum_{K/2<d\le K}d\varphi(d)+O(K^2)\\
 &=\frac7{2\pi^2}K^3+O(K^2\log K).
\end{align*}
Hence
\begin{equation}\label{eq:phi-cost}
 \sum_{d\le K}c_{K,d}\varphi(d)
 \le\frac{385}{64\pi^2}K^3+O(K^2\log K).
\end{equation}

Finally the cyclotomic product formula gives, uniformly in $d$,
\[
 \log\Phi_d(p)=\varphi(d)\log p
    +\sum_{a\mid d}\mu(d/a)\log(1-p^{-a})
    =\varphi(d)\log p+O_p(1).
\]
The uniform error follows from
$\sum_{a\ge1}|\log(1-p^{-a})|<\infty$.
Since $\sum_{d\le K}c_{K,d}=O(K^2\log K)$, summing these errors
together with~\eqref{eq:phi-cost} proves~\eqref{eq:denom-growth} and
completes Proposition~\ref{prop:denominators}.

\begin{corollary}\label{cor:integer-polys}
The polynomials
$P_K(X)=C_Kp^{-e_K}\Delta_K(X)$ belong to $\Z[X]$ and have degree $K$.
\end{corollary}
\begin{proof}
For primes $\ell\nmid p$, the specialization
of~\eqref{eq:laurent-integral} shows that $C_K\Delta_K$ is
$\ell$-integral, and division by $p^{e_K}$ preserves this property.
For primes $\ell\mid p$, Proposition~\ref{prop:content} gives the
integrality of $p^{-e_K}\Delta_K$ directly; multiplication by $C_K$
preserves it. A rational number integral at every prime is an integer.
The degree assertion follows from~\eqref{eq:leading}.
\end{proof}

\section{Infinitely many nonzero determinants}

This argument concerns the actual real functional. It requires no
rationality hypothesis on $\xi_p$ and no assertion about a sign pattern.

\begin{proposition}\label{prop:normality}
The determinants $\Delta_K(\xi_p)$ are nonzero for infinitely many
positive integers $K$.
\end{proposition}

\begin{proof}
Define the Cauchy transform, initially for $|z|>q$, by
\begin{equation}\label{eq:cauchy}
 \mathcal F(z)=\Lcal\left(\frac1{z-x}\right)
   =\sum_{m\ge1}\frac{q^m(z+q^m)}{(z-q^m)^3}.
\end{equation}
This series converges locally normally away from
$\{0,q,q^2,\ldots\}$. On a compact subset of that domain the tail
terms are uniformly $O(q^m)$; only finitely many earlier terms remain.
It is analytic at infinity and is $O(z^{-2})$. At $z=q^m$, its
principal part is
\[
 2q^{2m}(z-q^m)^{-3}+q^m(z-q^m)^{-2}.
\]
All other summands are analytic in a neighborhood of that point.
Thus it has infinitely many genuine poles and is not rational.

For a polynomial $Q$ of degree at most $n-1$, set
\begin{align}
 R_{n,Q}(z)&=\Lcal\left(\frac{Q(x)}{(z-x)D_n(x)}\right),
 \label{eq:Rnormal}\\
 T_{n,Q}(z)&=\Lcal\left(
 \frac{Q(z)D_n(x)-D_n(z)Q(x)}{(z-x)D_n(x)}\right).
 \label{eq:Tnormal}
\end{align}
The numerator in~\eqref{eq:Tnormal} vanishes identically on $z=x$,
so division by $z-x$ is exact as a polynomial in $x,z$. Applying
$\Lcal$ to the coefficients after dividing by $D_n(x)$ produces a
polynomial $T_{n,Q}(z)$ of degree at most $n-1$. Linearity gives
\begin{equation}\label{eq:interpolation-id}
 \mathcal F(z)Q(z)-T_{n,Q}(z)=D_n(z)R_{n,Q}(z).
\end{equation}

Suppose $\Delta_n(\xi_p)=0$. A nonzero vector in the kernel of its
Gram matrix supplies a nonzero polynomial $Q_n$ of degree at most
$n-1$ satisfying
\begin{equation}\label{eq:orthogonality}
 \Lcal\left(\frac{x^jQ_n(x)}{D_n(x)}\right)=0
       \qquad(0\le j<n).
\end{equation}
The geometric expansion of $(z-x)^{-1}$ at infinity, followed by
\eqref{eq:orthogonality}, gives $R_{n,Q_n}(z)=O(z^{-n-1})$.
All termwise operations are legitimate: for fixed $n$, the rational
functions are analytic near $[0,q]$, their differentiated summands
are $O(q^m)$, and the geometric expansion and its first two
$x$-derivatives converge uniformly for $|z|>r>q$.
Writing $T_n=T_{n,Q_n}$, we therefore have
\begin{equation}\label{eq:errornormal}
 E_n(z):=\mathcal F(z)Q_n(z)-T_n(z)=O(z^{-1}).
\end{equation}
The transform $R_{n,Q_n}$ is analytic at every $p^a$, all of which
lie outside the support $\{q^m:m\ge1\}$. Thus~\eqref{eq:interpolation-id}
also gives
\begin{equation}\label{eq:exterior-zeros}
 E_n(p^a)=0\qquad(1\le a\le n).
\end{equation}

If the determinants of both sizes $n$ and $n+1$ vanish, choose
$Q_n,Q_{n+1}$ as above. The polynomial
\[
 V(z)=T_{n+1}(z)Q_n(z)-T_n(z)Q_{n+1}(z)
     =Q_{n+1}(z)E_n(z)-Q_n(z)E_{n+1}(z)
\]
is $O(z^{n-1})$ at infinity, using the degree bounds on the $Q$'s
and~\eqref{eq:errornormal}. Hence it has degree at most $n-1$.
It vanishes at the $n$ distinct points $p,p^2,\ldots,p^n$ by
\eqref{eq:exterior-zeros}, so it is zero. Consequently
$T_n/Q_n=T_{n+1}/Q_{n+1}$ as rational functions. This argument never
divides by $Q_n(p^a)$ or $Q_{n+1}(p^a)$; coincident zeros do not
invalidate it.

Assume now, toward contradiction, that all determinants from some
size $N$ onward vanish. Their rational functions are all the same;
choose a fixed reduced expression $T/Q$ with $Q\ne0$ and
$\gcd(T,Q)=1$. This includes the possibility $T=0$, with $Q=1$.
Polynomial coprimality yields nonzero polynomials $U_n$ such that
\[
 Q_n=U_nQ,\qquad T_n=U_nT\qquad(n\ge N).
\]
Set $E=\mathcal FQ-T$. Since $\mathcal F$ is nonrational, $E$ is
nonzero. It is meromorphic at infinity, and
$U_NE=E_N=O(z^{-1})$ implies a finite positive order of vanishing
there. Thus, for some integer $h\ge1$ and $c\ne0$,
\[
 E(z)=cz^{-h}+O(z^{-h-1}).
\]
The identities $U_nE=E_n=O(z^{-1})$ force $\deg U_n\le h-1$ for
every $n\ge N$. The same Laurent expansion implies that $E(z)$ has
no zeros when $|z|$ is sufficiently large. For all such exterior
nodes $p^a$ with $a\le n$, equation~\eqref{eq:exterior-zeros} forces
$U_n(p^a)=0$. As $n$ tends to infinity, this gives arbitrarily many
distinct zeros of a nonzero polynomial of degree at most $h-1$, a
contradiction. There can be no eventually singular tail.
\end{proof}

\section{Assembly of the rationality contradiction}

\begin{proof}[Proof of the candidate main theorem]
Let
\[
 \eps=\frac{11}{18}-\frac{385}{64\pi^2}.
\]
This number is strictly positive. Indeed, the equivalent inequality
is $\pi^2>315/32$, which follows already from $\pi>157/50$.
For an entirely rational lower estimate,
\[
 \eps>
 \frac{11}{18}-\frac{385}{64(157/50)^2}
 =\frac{3487}{3549456}>0.
\]

Corollary~\ref{cor:integer-polys} supplies degree-$K$ integer
polynomials $P_K=C_Kp^{-e_K}\Delta_K$.
Combining Proposition~\ref{prop:analytic},
Proposition~\ref{prop:denominators}, and $e_K=K^3/2-K/2$, gives
\begin{align*}
 \log|P_K(\xi_p)|
 &\le\frac{385}{64\pi^2}(\log p)K^3
       -\frac12(\log p)K^3-\frac19(\log p)K^3
       +O_p(K^2\log K)\\
 &=-\eps(\log p)K^3+O_p(K^2\log K).
\end{align*}
The inequality is also harmless at zero values. By
Proposition~\ref{prop:normality}, infinitely many values are nonzero.

If $\xi_p=a/b$ with $a\in\Z$ and $b\in\Z_{>0}$, then for each $K$
the number $b^KP_K(a/b)$ is an integer, because $P_K\in\Z[X]$ and
$\deg P_K=K$. On the infinite set of nonzero values it satisfies
\[
 1\le |b^KP_K(a/b)|
 \le\exp\{-\eps(\log p)K^3+O_p(K^2\log K)+K\log b\}.
\]
The upper bound tends to zero as $K\to\infty$. Every infinite set
of positive integers is unbounded, so this is a contradiction along
the stated infinite set. Hence $\xi_p$ is irrational.
\end{proof}

\begin{remark}[Scope of the conclusion]
The argument, if validated, proves irrationality for each fixed
integer $p\ge2$. It does not supply a uniform effective irrationality
measure, transcendence, a statement for arbitrary rational $q$, or
a result about the ordinary real number $\zeta(7)$. No such extension
is inferred here. The nonvanishing proposition needs only infinitely
many sizes, and neither the size of the first useful determinant nor
the density of those sizes is asserted.
\end{remark}

\begin{remark}[What the arithmetic proof includes]
The cyclotomic constant in this manuscript is the conservative
$385/(64\pi^2)$. The proof explicitly retains the moment
$d^2/(z^d-1)$ in the range $K/2<d\le K$ and bounds its exact rank
by $\max(0,K-d-1)$. The stronger variant proved in the appendix
is not needed for the main argument. The real determinant estimate
allows vanishing; nonvanishing is supplied separately by
Proposition~\ref{prop:normality}.
\end{remark}

\appendix
\section{A stronger high-order denominator bound}\label{app:stronger}

The main proof uses weights $(2,1)$ on double clusters. A basis that
vanishes exactly at the upper node permits weights $(2,0)$ and a
larger asymptotic margin. This appendix records the separate improvement;
none of its assertions is needed to make the main coefficient
$\eps$ positive.

\begin{proposition}\label{prop:stronger}
In~\eqref{eq:cd}, one may replace $7(K-d)+3$ by $5(K-d)+3$ for
$K/2<d\le K$. If $C_K^{\mathrm{str}}$ is the resulting cyclotomic
product and $P_K^{\mathrm{str}}=C_K^{\mathrm{str}}p^{-e_K}\Delta_K$,
then $P_K^{\mathrm{str}}\in\Z[X]$ and
\begin{align*}
 \log C_K^{\mathrm{str}}
 &\le\frac{321}{64\pi^2}(\log p)K^3+O_p(K^2\log K),\\
 \log|P_K^{\mathrm{str}}(\xi_p)|
 &\le-\left(\frac{11}{18}-\frac{321}{64\pi^2}\right)
            (\log p)K^3+O_p(K^2\log K).
\end{align*}
\end{proposition}

\begin{proof}
Retain the local notation of Section~4, with $d>K/2$ and $t=K-d$.
The double clusters are exactly $\{a,a+d\}$ for $1\le a\le t<d$.
For one such cluster let
\[
 F_a(x)=(x-z^a)(x-z^{a+d}),\qquad V_a(x)=F(x)/F_a(x),
\]
and replace its two basis elements by
\[
 \widetilde G_{a,0}(x)=V_a(x),\qquad
 \widetilde G_{a,1}(x)=V_a(x)(x-z^{a+d}).
\]
The change from the previous cluster basis is triangular with diagonal
entries one: $x-z^{a+d}$ differs from $x-\omega^a$ by an analytic
scalar. Thus the total basis change remains a local unit. Retain the
singleton basis elements from~\eqref{eq:cluster-basis}.

The proper part is still block diagonal by clusters. In the block
for $a$, every entry involving $\widetilde G_{a,1}$ has zero residue
at the upper node $z^{a+d}$, because that basis element vanishes
there identically for variable $z$. At the lower node $z^a$, the
single factor $z^a-z^{a+d}$ in $\widetilde G_{a,1}(z^a)$ has local
valuation one and cancels the one vanishing factor in $F'(z^a)$.
Since $a<d$, the lower moment $X-S_a(z)$ is analytic.
Consequently, the proper-part entries $(0,1)$, $(1,0)$ and $(1,1)$
are analytic. The $(0,0)$ entry has pole order at most four: at
either node the simple-pole coefficient has order at most one,
and the pole moment has order at most three.

Weights $(2,0)$ therefore clear the proper part of this block at
determinant cost four. The singleton $a=d$ still costs three, and
all other singletons cost zero. The total cost is $4t+3$.
The polynomial part is treated exactly as before: its only possible
singular moment is $\mu_d$, and its coefficient form has exact rank
at most $\max(0,t-1)$ by~\eqref{eq:rank}. All new weights are
nonnegative, so the analytic coefficient matrices stay integral after
scaling. The same column-expansion argument gives
\[
 -v(\Delta_{K,z})\le4t+3+\max(0,t-1)\le5t+3.
\]
This proves the improved local bound and, by the same localization
argument, the improved integral-polynomial assertion.

The low-order contribution remains $161/(64\pi^2)$. Replacing the
high-order coefficient seven by five makes its contribution
$5/(2\pi^2)$. Their sum is $321/(64\pi^2)$, with the same
$O_p(K^2\log K)$ error as in Section~4. The real estimate and
the coefficientwise $p$-factor are unchanged, yielding the second
inequality of the proposition.
\end{proof}

The margin in this variant is
$11/18-321/(64\pi^2)$, approximately $0.1029220494$.
This decimal only describes the exact expression; it is not a
numerical input to the proof.

\begin{thebibliography}{9}

\bibitem{JM}
F.~Jouhet and E.~Mosaki,
\emph{Irrationalit\'e aux entiers impairs positifs d'un
$q$-analogue de la fonction z\^eta de Riemann},
International Journal of Number Theory (2010),
\href{https://doi.org/10.1142/S1793042110003332}{doi:10.1142/S1793042110003332};
\href{https://arxiv.org/abs/0712.1762}{arXiv:0712.1762}.

\bibitem{Zudilin2016}
W.~Zudilin,
\emph{On the irrationality of generalized $q$-logarithm},
Research in Number Theory \textbf{2}, article 15 (2016),
\href{https://doi.org/10.1007/s40993-016-0042-x}{doi:10.1007/s40993-016-0042-x};
\href{https://arxiv.org/abs/1601.02688}{arXiv:1601.02688}.

\bibitem{Tauraso}
R.~Tauraso,
\emph{A $q$-analogue of the Koecher--Leshchiner generating function
of odd zeta values}, preprint (2025),
\href{https://arxiv.org/abs/2511.18034}{arXiv:2511.18034}.

\bibitem{KY}
J.~Koizumi and A.~Yokoi,
\emph{Ap\'ery-type approximations and irrationality measures for
certain $q$-series}, preprint (2026),
\href{https://arxiv.org/abs/2608.26918}{arXiv:2608.26918}.

\end{thebibliography}
\end{document}
